\documentclass[10pt,letterpaper]{article}
\usepackage[T1]{fontenc}
\usepackage[utf8]{inputenc}
\usepackage{lmodern,microtype,amsmath,amssymb,amsthm,mathtools,booktabs,enumitem}
\usepackage[margin=0.9in]{geometry}
\usepackage[colorlinks=true,linkcolor=black,urlcolor=blue,citecolor=black]{hyperref}
\usepackage{fancyhdr}
\setlength{\headheight}{14pt}
\pagestyle{fancy}\fancyhf{}\lhead{From Prephysical Hypotheses to Physical Models}\rhead{Mathematical research supplement}\cfoot{\thepage}
\newtheorem{theorem}{Theorem}
\newtheorem{lemma}[theorem]{Lemma}
\newtheorem{corollary}[theorem]{Corollary}
\newtheorem{proposition}[theorem]{Proposition}
\newtheorem{remark}[theorem]{Remark}
\newcommand{\HH}{\mathcal H}
\newcommand{\DD}{\mathcal D}
\newcommand{\SSep}{\mathsf{Sep}}
\newcommand{\KK}{\mathcal K}
\newcommand{\MM}{\mathcal M}
\newcommand{\tr}{\operatorname{Tr}}
\newcommand{\id}{\mathbf 1}
\newcommand{\C}{\mathbb C}
\newcommand{\LOCC}{\mathrm{LOCC}}
\newcommand{\GHZ}{\mathrm{GHZ}}
\title{\textbf{Instrument Closure and Multipartite Selection:}\\[4pt]
Exact Counterexamples and Conditional Completion Theorems}
\author{Kirtus Justus\\[-1pt]\small Research direction and scholarly responsibility\\
\small Port Coquitlam, British Columbia, Canada}
\date{12 November 2026, B.C. Pacific Time\\[3pt]
\small Initial External Grounding; mathematics under stipulated premises; no ontology selected}
\begin{document}
\maketitle
\begin{abstract}
A two-system convex family can reproduce ordinary elementary complex two-level
quantum statistics while restricting composite preparations. An earlier four-system
counterexample showed that a uniform convex interpolation between fully separable
and all density operators fails parallel-preparation closure. We strengthen this in
four ways. First, every nonendpoint interpolation already fails closure under a
\emph{single-party successful local filter}, and the violation and its nonzero
success probability are computed exactly. Second, permitting a different mixture
parameter at four systems cannot restore \emph{both} parallel preparation and
partial trace. Third, an explicit matching-generated hierarchy provides a positive
countercontrol: it is closed under parallel preparation, discarding and elementary
local trace-preserving operations, retains entangled two-system preparations and
may violate CHSH, but excludes a three-system GHZ state. Fourth, within a stipulated
complex density-operator representation, closure under local successful filters
and local unitaries makes \emph{one entangled pure two-system seed} sufficient to
fill the complete two-system state space. A related mixed-state extension invokes
published distillation results and substantially stronger assumptions. Every
statement is conditional mathematical research, not an emergence of operational
probability, physical systems, time, quantum mechanics, classical mechanics, or a
calibrated prephysical-to-physical map.
\end{abstract}

\section{Assessment contract and definitions}
The program's target is an explained emergence of quantum and classical mechanics
from independently warranted prephysical structure, with an empirical map.
Here the \emph{mathematical control} is assumed: an elementary Hilbert space
$H=\C^2$ with its usual complex inner product, positive density operators,
Born-form evaluation, Hilbert tensor products and the standard meaning of a
quantum operation. These are \emph{stipulations of the test}, not derived
primitive phenomena and not identification of physical qubits.
Let $\DD_n$ be the density operators on $H^{\otimes n}$, and
\[
\SSep_n=\operatorname{conv}\{\rho_1\otimes\cdots\otimes\rho_n:
\rho_j\in\DD_1\}
\]
the fully separable subset. For $0\le t\le1$, define the Minkowski mixture
\begin{equation}\label{eq:K}
\KK_t^{(2)}=(1-t)\SSep_2+t\DD_2
=\{(1-t)\sigma+t\tau:\sigma\in\SSep_2,\ \tau\in\DD_2\}.
\end{equation}
It is closed under convex mixing and elementary local completely positive
trace-preserving (CPTP) operations. Its extreme $t=0$ is a \emph{separable
composite with quantum elementary systems}, not classical Hamiltonian mechanics.
Its $t=1$ endpoint is the ordinary quantum two-system composite.
Set
\[
|\Phi\rangle=\frac{|00\rangle+|11\rangle}{\sqrt2},\qquad
P=|\Phi\rangle\langle\Phi|,\qquad
\rho_t=(1-t)|00\rangle\langle00|+tP.
\]
We use $0\preceq E\preceq\id$ as positive operator inequalities. A linear
separation functional need not already be implemented as a physical effect.
When operational access is discussed, the needed admissible instrument is stated.

\begin{lemma}[Sharp Bell support]\label{lem:support}
For every $0\le t\le1$,
\[
\max_{\rho\in\KK_t^{(2)}}\tr(P\rho)=\frac{1+t}{2}.
\]
The state $\rho_t$ attains the maximum.
\end{lemma}
\begin{proof}
For product unit vectors $|a\rangle|b\rangle$, Cauchy--Schwarz gives
$|\langle\Phi|a\otimes b\rangle|^2\le1/2$; convexity extends the
bound to $\SSep_2$. The maximum over $\DD_2$ is $1$.
The affine mixture has upper bound $(1-t)/2+t$, attained by the
separable $|00\rangle\langle00|$ and the quantum $P$ constituents.
\end{proof}

\section{Instrument-level obstruction at two systems}
\begin{theorem}[A single successful local filter leaves every interior family]
\label{thm:filter}
For every $0<t<1$, $\KK_t^{(2)}$ is not closed under normalized
nonzero successful branches of all single-party local quantum instruments.
In fact, an ordinary two-outcome instrument on the first subsystem has a branch
with success probability exactly $t$, whose conditioned output is outside
$\KK_t^{(2)}$.
\end{theorem}
\begin{proof}
For $0<x\le1$ use Kraus operators on the first system
\[
F_x=\begin{pmatrix}x&0\\0&1\end{pmatrix},\qquad
F_x^{\rm fail}=\begin{pmatrix}\sqrt{1-x^2}&0\\0&0\end{pmatrix},\qquad
F_x^\dagger F_x+(F_x^{\rm fail})^\dagger F_x^{\rm fail}=\id_2.
\]
Applying the successful branch $F_x\otimes\id_2$ to $\rho_t$ gives the
unnormalized state
\[
\widetilde\rho_{t,x}=(1-t)x^2|00\rangle\langle00|
+\frac t2(x|00\rangle+|11\rangle)(x\langle00|+\langle11|),
\]
with success probability
\begin{equation}\label{eq:success}
p_{t,x}=\tr\widetilde\rho_{t,x}=
\frac{(2-t)x^2+t}{2}>0.
\end{equation}
The conditioned Bell-projector expectation is
\begin{equation}\label{eq:fidelity}
f_{t,x}=\frac{\tr(P\widetilde\rho_{t,x})}{p_{t,x}}
=\frac12+\frac{tx}{(2-t)x^2+t}.
\end{equation}
The derivative vanishes at $x_*=\sqrt{t/(2-t)}\in(0,1)$,
where it is maximal on $(0,1]$. At this point
\begin{equation}\label{eq:gap}
p_{t,x_*}=t,\qquad
f_{t,x_*}=\frac{1+\sqrt{t/(2-t)}}{2},\qquad
f_{t,x_*}-\max_{\rho\in\KK_t^{(2)}}\tr(P\rho)
=\frac{\sqrt{t/(2-t)}-t}{2}>0.
\end{equation}
Strict positivity follows from
$\frac{t}{2-t}-t^2=\frac{t(1-t)^2}{2-t}>0$.
Lemma~\ref{lem:support} therefore excludes the normalized output from
$\KK_t^{(2)}$. This is a violation of \emph{postselected-state closure},
not of trace-preserving channel closure.\end{proof}

\begin{remark}[Access, limits and controls]
At $t=1/2$, $x_*=1/\sqrt3$, success probability is $1/2$,
$f_{t,x_*}=(1+1/\sqrt3)/2$, and the gap is $0.03867513459\ldots$.
As $t\downarrow0$ success probability tends to zero; the endpoint
$t=0$ is not assigned a nonexistent normalized zero-probability branch.
Both $\KK_0^{(2)}=\SSep_2$ and $\KK_1^{(2)}=\DD_2$ are closed under
nonzero-probability local filters. Near $t\uparrow1$ the positive gap is
$\frac14(1-t)^2+O((1-t)^3)$; floating-point zero is not a mathematical
counterexample. The Bell witness is accessible by \emph{product-observable
statistics} if local Pauli measurements are supplied, since
$P=\tfrac14(\id\otimes\id+X\otimes X-Y\otimes Y+Z\otimes Z)$.
Those effects, outcome probabilities, postselection, Pauli measurements and
preparation resources are all imported in this mathematical test.
\end{remark}

\section{An arity-dependent interpolation still fails at four systems}
Allow two and four elementary systems to use \emph{different} mixture weights
$t,s\in[0,1]$:
\[
\KK_t^{(2)}=(1-t)\SSep_2+t\DD_2,\qquad
\KK_s^{(4)}=(1-s)\SSep_4+s\DD_4.
\]
\begin{theorem}[Independent preparation and discarding are incompatible]
\label{thm:arity}
If (i) the tensor product of any two allowed pair states is admitted at four
systems and (ii) discarding either disjoint pair maps every four-system
state into $\KK_t^{(2)}$, then necessarily $t=s\in\{0,1\}$.
In particular no interior pair weight can be repaired by choosing an
independent four-body weight $s$ of the displayed form.
\end{theorem}
\begin{proof}
\emph{Discard bound.} For any $s$, the state
$\rho_s^{12}\otimes|00\rangle\langle00|^{34}$ belongs to $\KK_s^{(4)}$:
its fully separable and unrestricted density constituents can be chosen
with the same spectator $|00\rangle^{34}$. Discarding systems 3,4 gives
$\rho_s^{12}$; Lemma~\ref{lem:support} implies
\begin{equation}\label{eq:up}s\le t.\end{equation}
\emph{Parallel-preparation bound.} Let $P_{12}=P\otimes\id_4$ and
$P_{34}=\id_4\otimes P$ and define the projector
\[
E=\id_{16}-(\id_4-P)\otimes(\id_4-P).
\]
For any fully product pure four-party state, the two disjoint pair Bell
probabilities $f_1,f_2$ satisfy $f_i\le1/2$. Thus
$\tr(E\omega)=1-(1-f_1)(1-f_2)\le3/4$ for all
$\omega\in\SSep_4$. The upper bound is attained by $|0000\rangle$;
$\max_{\DD_4}\tr E=1$. Consequently
\[
\max_{\omega\in\KK_s^{(4)}}\tr(E\omega)=\frac34+\frac s4.
\]
Independent preparation of $\rho_t^{12}$ and $\rho_t^{34}$ requires
$\rho_t\otimes\rho_t\in\KK_s^{(4)}$, whose actual witness value is
$1-((1-t)/2)^2$. Hence
\begin{equation}\label{eq:down}
s\ge2t-t^2.
\end{equation}
Together \eqref{eq:up} and \eqref{eq:down} require
$t\ge2t-t^2$, i.e. $t(1-t)\le0$. With $t\in[0,1]$, this allows only
$t=0$ or $1$, and the two inequalities then force $s=t$.
\end{proof}
\begin{remark}
The theorem excludes precisely a \emph{uniform-at-each-arity
Minkowski interpolation} with weights allowed to vary by arity. It does
not exclude all possible four-system extensions. The witness $E$ can
serve as an ordinary linear separation functional even if a proposed
instrument theory does not physically implement it. For the special case
$s=t$, the violation is $t(1-t)/4$, recovering the earlier four-party
counterexample. The new contradiction uses an additional \emph{discarding}
obligation and defeats attempts to change $s$ to repair that example.
\end{remark}

\section{Positive countercontrol: a coherent restricted hierarchy}
A no-go must not be promoted beyond its hypotheses. Let $\mathfrak P_n$
be the finite set of partitions of $\{1,\ldots,n\}$ into blocks of size
one or two. Set $L_1=\DD_1$, $L_2=\KK_t^{(2)}$ and define
\begin{equation}\label{eq:matching}
\MM_n(t)=\operatorname{conv}\bigcup_{\pi\in\mathfrak P_n}
\left\{\bigotimes_{b\in\pi}\rho_b:\rho_b\in L_{|b|}\right\}.
\end{equation}
The tensor factors are placed on their labeled subsystems. This is a
\emph{matching-generated} family, not the interpolation of
Theorem~\ref{thm:arity}.
\begin{theorem}[Consistency under specified weak operations]
\label{thm:matching}
For every $t\in[0,1]$, the hierarchy \eqref{eq:matching} is compact,
convex, invariant under permutations and closed under independent
preparation, discarding any subset of elementary systems, and arbitrary
\emph{unconditioned} one-site CPTP operations. It also permits destructive
single-site measurements with postselection \emph{if the measured site is
removed}. Its single-system and two-system spaces are $\DD_1$ and
$\KK_t^{(2)}$. Product-observable tomography distinguishes its density
operators. Nevertheless $\MM_3(t)\ne\DD_3$ for every $t$.
\end{theorem}
\begin{proof}
Each block-state set is compact, and there are finitely many partitions;
finite-dimensional convex hull preserves compactness. Permutations relabel
partitions. The tensor product of two block-product terms is a block-product
term on the disjoint union; bilinear expansion extends this to convex
mixtures. Partial traces either remove whole blocks or turn a surviving
member of a two-block into an arbitrary $\DD_1$ state. On-site CPTP
operations preserve $L_1$ and preserve $L_2$ because they separately
preserve the fully separable and quantum constituents. A destructive
single-site effect similarly leaves its block partner in $\DD_1$ and
other blocks unchanged, with normalized mixtures remaining admissible.
Product Pauli effects span the full Hermitian operator space and hence
separate any two different states.

Let $|\GHZ\rangle=(|000\rangle+|111\rangle)/\sqrt2$ and
$Q=|\GHZ\rangle\langle\GHZ|$. Every block-product term at three
systems factors across some one-site versus two-site bipartition. Across
every such split the Schmidt coefficients of $|\GHZ\rangle$ are
$(1/\sqrt2,1/\sqrt2)$. Thus $\tr(Q\omega)\le1/2$ for every
$\omega\in\MM_3(t)$, by the maximal Schmidt-overlap bound and convexity.
But $\tr(Q Q)=1$, so $Q\notin\MM_3(t)$.
\end{proof}
\begin{remark}[Not a full operational emergence theory]
The hierarchy deliberately does \emph{not} claim closure under
nondestructive successful local filters, arbitrary joint instruments,
entangling gates or joint operations on all members of one laboratory.
Adding these assumptions requires a new closure proof and can defeat
\eqref{eq:matching}. For every $t>0$ it nevertheless contains the
entangled $\rho_t$ (its partial transpose has eigenvalue $-t/2$).
For the stipulated Pauli CHSH settings it contains a state with
$\langle\mathrm{CHSH}\rangle=\sqrt2(1+t)>2$ whenever
$t>\sqrt2-1$. Thus all-arity parallel preparation, discard, local
state tomography, and a Bell-CHSH violation do not alone force the
unrestricted many-body density spaces. This is a restricted mathematical
quantum subtheory with imported probabilities and operations, \emph{not}
a proposed physical ontology.
\end{remark}

\section{A conditional positive completion mechanism}
\begin{theorem}[One entangled pure seed with local selective filters]
\label{thm:pure}
Let $\mathcal L\subseteq\DD_2$ be convex and invariant under all local
unitaries and every normalized nonzero successful branch of a one-party
local quantum filter $F\otimes\id_2$, where $F^\dagger F\preceq\id_2$.
If $\mathcal L$ contains even one \emph{pure entangled} two-system state,
then $\mathcal L=\DD_2$.
\end{theorem}
\begin{proof}
By local unitaries write the seed in Schmidt form
$|\psi\rangle=a|00\rangle+b|11\rangle$, with $a,b>0$ and
$a^2+b^2=1$. Every target two-system unit vector has Schmidt form
$(U\otimes V)(c|00\rangle+d|11\rangle)$ with $c,d\ge0$,
$c^2+d^2=1$. Choose any $\kappa>0$ small enough that
$\kappa c/a\le1$ and $\kappa d/b\le1$. Then
\[
F=\operatorname{diag}(\kappa c/a,\kappa d/b)
\quad\text{satisfies}\quad F^\dagger F\preceq\id_2,
\]
and $(F\otimes\id)|\psi\rangle=\kappa(c|00\rangle+d|11\rangle)$.
The branch has positive success probability $\kappa^2$; normalization,
followed by the local $U\otimes V$, produces the target pure state.
Thus $\mathcal L$ contains every pure bipartite state. Their convex hull
is $\DD_2$, while $\mathcal L\subseteq\DD_2$ by hypothesis.
\end{proof}
\begin{corollary}[Mixed entangled seed under stronger published premises]
\label{cor:mixed}
Suppose $\mathcal L\subseteq\DD_2$ is also compact and is the two-system
sector of an all-arity theory closed under tensor powers, discarding,
normalized successful finite-round $\LOCC$ protocols on arbitrarily many
copies, local unitaries and successful local filters. If $\mathcal L$
contains an entangled two-system density operator, then
$\mathcal L=\DD_2$, \emph{conditional on the published theorem} that
every entangled $2\times2$ quantum state is distillable by such protocols.
\end{corollary}
\begin{proof}
The two-qubit distillation theorem of Horodecki, Horodecki and
Horodecki~\cite{Horodecki1997} gives finite-copy $\LOCC$ outputs with
Bell-state fidelity tending to one. The assumed operational closure
places every normalized output in $\mathcal L$. Compactness places
the limiting Bell projector in $\mathcal L$. Apply
Theorem~\ref{thm:pure}.\end{proof}
\begin{remark}[Assumption debt]
Theorem~\ref{thm:pure} assumes a standard complex composite, a pure
entangled seed, arbitrarily chosen local filters, successful-branch
updating, and convex mixing. Corollary~\ref{cor:mixed} adds all-arity
systems, cross-copy laboratory operations, classical communication,
postselection, a compactness limit, and an external distillation theorem.
These assumptions can do the mathematical selection work but do not explain
why a prephysical substrate supplies them. The statement is \emph{not}
that any CHSH violation, bare tensor closure, or holonomy alone implies
full quantum theory.
\end{remark}

\section{Reconciliation with the broader research objective}
The projective holonomy determinant identity, its root-of-unity ambiguity,
global conjugation equivalence, finite exact-homogeneity obstruction,
and finite-cycle complex-type algebra remain conditional or restricted
results of earlier work. Neither they nor the present composition theorems
derive an operational clock, the Born rule, a Hamiltonian, physical
spacetime, a quantum-to-classical Hamiltonian limit, or empirical
prediction with uncertainty. The finite-cycle construction also admits a
classical simplex completion; present composite results import complex
Hilbert spaces and therefore \emph{do not bridge that selection gap}.

The resulting obligation hierarchy is sharper: choose independently
warranted primitive operations and records; identify probability and
measurement update; justify composition and conditioning at every arity;
identify any actual entangling preparation and instrument; then separately
derive dynamics, physical units and a controlled classical limit. A
failure of one natural interpolation does not exclude every completion;
Theorem~\ref{thm:matching} supplies an explicit countercontrol. Work
remains at Initial External Grounding. No ontology is selected.

\section*{Reproducibility and provenance}
The accompanying source checker uses exact SymPy reductions, high-precision
MPMath boundary cases, ordinary finite-dimensional matrix controls, a
seeded batch of GHZ biseparable controls and sampled pure-state local
filters. Theorems, rather than finite random tests, supply the universal
quantifiers. The one-party local-filter result and arity-varying four-party
bound have exact algebraic proofs; the sampling checks do not estimate a
physical distribution. All boundary and numerical outlier cases are retained.

Research direction and scholarly responsibility: Kirtus Justus.
Mathematical development, code execution and internal checking assisted by
ChatGPT (GPT-6). This is one-assistant internal work and incorporates
other assistant-generated submissions as \emph{internal research provenance},
not as independent refereeing. The underlying research release was not
modified; no experimental results, independent scientific validation or
host-enforced research protocol are claimed.

\begin{thebibliography}{9}\small
\bibitem{Bennett1996} C. H. Bennett, H. J. Bernstein, S. Popescu and
B. Schumacher, ``Concentrating partial entanglement by local operations,''
\emph{Phys. Rev. A} \textbf{53}, 2046 (1996).
\url{https://doi.org/10.1103/PhysRevA.53.2046}.
\bibitem{Horodecki1997} M. Horodecki, P. Horodecki and R. Horodecki,
``Inseparable Two Spin-1/2 Density Matrices Can Be Distilled to a Singlet Form,''
\emph{Phys. Rev. Lett.} \textbf{78}, 574 (1997).
\url{https://doi.org/10.1103/PhysRevLett.78.574}.
\bibitem{Nielsen1999} M. A. Nielsen,
``Conditions for a Class of Entanglement Transformations,''
\emph{Phys. Rev. Lett.} \textbf{83}, 436 (1999).
\url{https://doi.org/10.1103/PhysRevLett.83.436}.
\bibitem{BarnumHilgert2020} H. Barnum and J. Hilgert,
``Spectral Properties of Convex Bodies,''
\emph{J. Lie Theory} \textbf{30}, 315--344 (2020).
\url{https://doi.org/10.5802/jolt.1118}.
\bibitem{Gelander2013} T. Gelander,
``Limits of Finite Homogeneous Metric Spaces,''
\emph{L'Enseignement Math\'ematique} \textbf{59}, 195--206 (2013).
\url{https://doi.org/10.4171/LEM/59-1-8}.
\end{thebibliography}
\end{document}
